\section{Introduction}\label{sec:intro}
Arbitrary-scale super-resolution (ASISR) reconstructs a high-resolution image from a low-resolution image and predicts its value at arbitrary coordinates. Local implicit decoders such as LIIF and LTE combine learned image features with coordinates and a target pixel size \citep{chen2021liif,lee2022lte}. A target scale changes both the frequencies directly observed in the input and the pixel footprint used for reconstruction. Aggregate PSNR can obscure which frequency bands improve, and a successful readout ablation does not explain why it improves them.

This paper asks two linked empirical questions: \emph{how does recovery change with frequency relative to the sampling scale, and where does a scale-dependent readout affect that recovery?} We make ASISR frequency recovery and paired readout behaviour the central objects of analysis. Spectral adaptation---changing a representation's preferred frequencies through Fourier features, activation scales or a covariance prior---provides context \citep{tancik2020fourier,sitzmann2020siren,kania2025fresh}. We do not propose another decoder or claim a general rule for selecting one.

Our measurement is the fraction of band energy recovered, relative to zero reconstruction of the centred signal. We compare trained networks with two specified stationary linear estimators and study frequency in low-resolution Nyquist units, $u=2s\nu$. A covariance-conditioned Wiener reference distinguishes an assumed prior from the empirical recovery of a trained network. It does not bound all possible estimators on a natural image. Classical mismatch costs explain tendencies in fixed linear channels; their transfer to band-averaged measurements and nonlinear networks is tested separately.

The strongest positive result concerns coarse structure. Frequency normalisation improves scale consistency in nine of nine model--dataset groups. Readout losses generally increase across the three frequency bands, and this ordering replicates in two additional training seeds for two architecture pairs. The stronger accounts fail. Normalisation does not meet the original collapse tolerance in three groups. A predictor based on linear recovery fails the original readout-loss correlation threshold. Scale marginalisation does not reproduce the ablated function, and its advantage over the simpler footprint predictor is architecture-dependent. We preserve these failures rather than replacing the hypotheses after seeing results.

\paragraph{Contributions and scope.}
\begin{enumerate}
\item An ASISR frequency-recovery baseline that separates covariance-conditioned linear references from measured network recovery, and records scale dependence across Set14, BSD100 and Urban100.
\item A replicated paired readout analysis: coarse band ordering is supported in four architecture--seed groups, whereas stronger magnitude, function-equivalence and mechanism claims are evaluated with their counterexamples.
\item An evidence boundary for analytical guidance: classical scalar and matrix identities, independently checked linear cost experiments, and a prospective negative comparison against cross-validation delimit what the ASISR diagnostics do and do not establish.
\end{enumerate}
The formulas are classical tools \citep{kay1993estimation,rasmussen2006gp}; the contribution claimed here is the organised empirical evidence and its scope. We distinguish original observations, pre-specified tests, independent reimplementations, saved-record audits and exploratory development. Extended protocols and results accompany the main text in the supplement.
